% ============================================================================
% October 13, 2026 build (agent draft). Appendix at the 7.9 production point.
% Additional policy arms: 7p9167_v6/table.md (read Oct 13). Arms marked
% PROVISIONAL there (stopped at the 30-map cap) are flagged in the table.
% Transfer arms: driver_policy_v6.py lines 1-5 and 195-210 (INS_KIND
% contingent / flat; 0.25 x shortfall below the median earnings state at each
% working age; unfunded, in the household budget only).
% The Oct 9 appendix on supply arms (14.402 figures) is retired; it is in git.
% ============================================================================
\input{sections/app_household_conditions}
\input{sections/app_model_figures}
% Appendix D, robustness of Section 6 (Oct 24, 2026). D.1: free entry, eta 1.5 (200-period arms, data_fi_torch/P15fi_*).
% D.2: slow stock (KMV permits), eta 0.63, 200-period fit (data_k200_torch/P063k_*). Numbers: output/model/experiments/
% housing_supply/20261021_paper_results/appD_values.json (p6_appendix.py); long run = stationary endpoint. At eta 1.5 the +10%
% supply run already matches credit on impact (1.585% vs 1.586%), so it is the matched experiment. Slow-stock matched supply (x1.4027,
% SUPMC) and credit for parents filled Oct 24 from data_k200_torch. Slow stock at eta 1.5 (P15k3_*, 300 periods) filled Oct 25 from
% appD_k15_values.json (appD_k15.py); matched = SUPM2 (+60%, births 2023 1.88% vs credit 1.92%).
% The previous version (old 100-period f3 runs) is in git (59d4c54e).
\providecommand{\pend}{\textcolor{red}{[pending]}}
\providecommand{\PRF}{../../output/model/experiments/housing_supply/20261021_paper_results}

\section{Robustness of the Policy Results}
\label{app:policy-robust}

This appendix repeats the experiments of Section~\ref{sec:policy} with a more
elastic supply and with a slowly adjusting stock.

\subsection{A more elastic supply}
\label{app:policy-eta}

Table~\ref{tab:app-policy-eta} repeats the experiments with free entry and
$\eta=1.5$. With a more elastic supply, a 10\% expansion lowers the price by
5.0\% on impact, against 8.6\% with $\eta=0.63$, and raises births in 2023 by
1.6\%, as much as the credit relaxation. It is therefore the matched
expansion. Credit again
moves births earlier: the completed fertility of the 2023 cohort falls by
1.0\%, while the supply expansion raises it by 2.6\%. In the steady state the price
is again the one at which fertility is at replacement, so credit for parents
raises it by about as much as with $\eta=0.63$, 34\%. With a more elastic
supply, more households fit at that price: 76\% more, against 36\%.

\begin{table}[H]
\centering
\caption{Counterfactuals from 2023 with $\eta=1.5$}
\label{tab:app-policy-eta}
\small
\begin{tabularx}{\textwidth}{@{}Xrrrcrr@{}}
\toprule
 & \multicolumn{3}{c}{Births} & Completed fertility & \multicolumn{2}{c}{House price} \\
\cmidrule(lr){2-4}\cmidrule(l){6-7}
 & 2023 & 2063 & Steady state & 2023 cohort & 2023 & Steady state \\
\midrule
Credit relaxation & $+1.6$ & $+0.1$ & $+4.8$ & $-1.0$ & $-0.1$ & $+2.6$ \\
Housing supply, $+10\%$ & $+1.6$ & $+4.3$ & $+10.0$ & $+2.6$ & $-5.0$ & $0.0$ \\
Credit for parents & $+5.1$ & $+14.2$ & $+75.6$ & $+5.8$ & $+0.1$ & $+33.9$ \\
\bottomrule
\end{tabularx}
\par\smallskip
\begin{minipage}{0.95\linewidth}
\footnotesize\textit{Notes:} As in Table~\ref{tab:policy-main}, with free
entry and $\eta=1.5$. The 10\% supply expansion raises births in 2023 as much
as the credit relaxation. Without policy, period fertility is 1.63 in 2023 and 1.67
in 2063.
\end{minipage}
\end{table}

\subsection{A slowly adjusting stock}
\label{app:policy-slow}

Here the housing stock adjusts slowly, as in
\citet{kaplanMitmanViolanteHousing2020}. Developers combine labor with building permits, which the government issues
in a fixed number each period. Construction and the stock then follow
\begin{equation}
 I_t=\delta_H H_0P_t^{\eta},\qquad H_t=(1-\delta_H)H_{t-1}+I_t .
 \label{eq:app-slow-stock}
\end{equation}
In a stationary equilibrium construction replaces depreciation, so the stock
lies on the same curve $H=H_0P^{\eta}$ as under free entry, and the 2007
equilibrium is the same. Along a transition only construction responds to the
price within a period, so the stock responds with an elasticity of about
$\delta_H\eta$, and the price absorbs most of a change in demand. The decline
in the preference for children from 2007 to 2023 is fitted under this rule, so
the 2023 state differs from the one in Section~\ref{sec:policy}. Everything
else is as in Section~\ref{sec:policy-design}.

Table~\ref{tab:app-policy-slow} reports the results with $\eta=0.63$. The credit
relaxation raises births by 1.7\% in 2023, but the gain falls to 0.6\% by 2063,
and the completed fertility of the 2023 cohort falls by 0.2\%. A 10\% supply
expansion lowers the price by 3.0\% in 2023 and raises births by only 0.4\%.
To raise births in 2023 as much as credit, the supply curve has to shift out
by 40\%. This matched expansion lowers the price by 10.9\% in 2023.
Construction keeps adding space, births are 12.2\% higher by 2063, and the
completed fertility of the 2023 cohort rises by 5.3\%. In the steady state the
price is back at its level without policy, so the supply expansions add as
many households as they add space, 10\% and 40\%, and credit adds 3.2\%.
Credit for parents raises births by 6.2\% in 2023 and by 32.9\% in the steady
state, when the price is 31.4\% higher.
Figure~\ref{fig:app-policy-slow-paths} shows the paths.

With $\eta=1.5$ the pattern is the same (lower panel of
Table~\ref{tab:app-policy-slow}). The credit relaxation raises births by 1.9\%
in 2023 and by 0.6\% in 2063, and the completed fertility of the 2023 cohort
falls by 0.8\%. To raise births in 2023 as much, the supply curve has to shift
out by about 60\%; this lowers the price by 10.8\% in 2023, raises births by
12.1\% by 2063 and the completed fertility of the 2023 cohort by 5.5\%.
Credit for parents raises births by 6.4\% in 2023 and the completed fertility
of the 2023 cohort by 7.5\%.

\begin{table}[H]
\centering
\caption{Counterfactuals from 2023 with a slowly adjusting stock}
\label{tab:app-policy-slow}
\small
\begin{tabularx}{\textwidth}{@{}Xrrrcrr@{}}
\toprule
 & \multicolumn{3}{c}{Births} & Completed fertility & \multicolumn{2}{c}{House price} \\
\cmidrule(lr){2-4}\cmidrule(l){6-7}
 & 2023 & 2063 & Steady state & 2023 cohort & 2023 & Steady state \\
\midrule
Credit relaxation & $+1.7$ & $+0.6$ & $+3.2$ & $-0.2$ & $+0.2$ & $+3.5$ \\
Housing supply, $+10\%$ & $+0.4$ & $+3.2$ & $+10.0$ & $+1.4$ & $-3.0$ & $0.0$ \\
Matched supply & $+1.7$ & $+12.2$ & $+40.3$ & $+5.3$ & $-10.9$ & $0.0$ \\
Credit for parents & $+6.2$ & $+17.1$ & $+32.9$ & $+5.9$ & $+1.8$ & $+31.4$ \\
\addlinespace
\multicolumn{7}{@{}l}{\emph{$\eta=1.5$}}\\
Credit relaxation & $+1.9$ & $+0.6$ & $+6.1$ & $-0.8$ & $-0.1$ & $+3.3$ \\
Housing supply, $+10\%$ & $+0.3$ & $+2.4$ & $+10.0$ & $+1.1$ & $-2.1$ & $0.0$ \\
Matched supply, $+60\%$ & $+1.9$ & $+12.1$ & $+60.0$ & $+5.5$ & $-10.8$ & $0.0$ \\
Credit for parents & $+6.4$ & $+17.2$ & $+64.7$ & $+7.5$ & $+1.1$ & $+29.8$ \\
\bottomrule
\end{tabularx}
\par\smallskip
\begin{minipage}{0.95\linewidth}
\footnotesize\textit{Notes:} As in Table~\ref{tab:policy-main}, $\eta=0.63$
(upper panel) and $\eta=1.5$ (lower panel), with the stock following
\eqref{eq:app-slow-stock}. Matched supply: the
supply shift that raises births in 2023 as much as the credit relaxation. Without
policy, period fertility is 1.64 in 2023 and 1.80 in 2063 with $\eta=0.63$, and
1.63 and 1.75 with $\eta=1.5$.
\end{minipage}
\end{table}

\begin{figure}[H]
\centering
\includegraphics[width=0.9\linewidth]{\PRF/fig_appD_paths_k200_eta063.pdf}
\caption{The credit relaxation and housing supply with a slowly adjusting stock}
\label{fig:app-policy-slow-paths}
\par\smallskip
\begin{minipage}{0.92\linewidth}
\footnotesize\textit{Notes:} \% differences from the no-policy path from the
same 2023 state, $\eta=0.63$, stock following \eqref{eq:app-slow-stock}.
Births in each four-year period (panel a), completed fertility of the cohort
aged 18--21 in the year shown (panel b), number of households (panel c) and
house price (panel d). Housing supply: the supply curve shifted out by 40\%,
which raises births in 2023 as much as the credit relaxation.
\end{minipage}
\end{figure}


\section{Additional Policy Results}
\label{app:policy-more}

Table~\ref{tab:app-policy-more} reports the policy experiments that the main
text does not discuss, with the same design as Section~\ref{sec:policy}: every
arm starts from the 2023 state and is compared with the no-policy path at the
same earnings risk and under the same housing supply.
\problem{Only the three base-risk rows at the top (property tax, insurance, flat
transfer) are at the constant-decay state. The fixed-stock, elasticity and
lower-risk rows are still from the earlier two-step path and are to be rerun.}

\paragraph{A rebated property tax.} The annual property tax rises from 1.06 to
2 percent of value from 2023 on, and the revenue is rebated equally to all
households. The tax lowers the price by 6.2 percent on impact, close to the
8.9 percent of the supply expansion, but births do not move ($-0.2$ percent)
and completed fertility rises by only 0.8 percent. The comparison is the
cleanest test of the mechanism of Section~\ref{sec:results}. The tax lowers
the price but raises the user-cost rate from 18.0 to 21.8 percent of the price
per period, so the per-period cost of a room rises; and the rebate returns the
revenue per household, not per room, so a family that needs more space pays
more for it. In this experiment, the lower purchase price is therefore not
equivalent to cheaper housing services. The comparison illustrates why the
price of a dwelling alone is insufficient to assess the fertility incentive.
\problem{The 13 percent increase in rent per room is a stationary
comparison; the response along the transition is [X.X].}

\paragraph{Earnings insurance and a flat transfer.} The insurance arm pays
every working household a quarter of the shortfall of its earnings below the
median earnings state at its age. The flat arm pays the expected value of that
transfer at each age to every working household. Neither is financed within
the model. Insurance raises births by 14.8 percent on impact and completed
fertility by 13.8 percent; the flat transfer by 8.0 and 7.8 percent. The
difference, about 6 percentage points, is the value of insuring the downside
of earnings rather than paying the same amount to everyone, as \citet{fraserIncomeRisk2001}
predicts for an irreversible commitment. Under lower earnings risk the
difference shrinks to 1.5 points on impact. Because neither transfer is
financed, they do not establish the effects of a budget-neutral program.
Financing could change both the distribution of resources and the households
whose fertility responds. The comparison isolates the difference between
state-contingent and age-specific transfers within the stated experiment;
it does not provide a fiscal or welfare ranking of implementable policies.

\paragraph{A fixed housing stock.} Holding the stock at its 2023 level,
instead of letting it follow the supply curve, lowers the price by 9.2 percent
in 2023 and 29 percent by 2063, as the population shrinks into the same stock;
completed fertility rises by 3.5 percent for the 2023 cohort and 8.6 percent
for the 2043 cohort. With the stock fixed, demand policies raise the price
more and births somewhat less. The credit relaxation raises the price by 2.1
percent on impact and leaves completed fertility unchanged ($+0.1$ percent),
the case in which credit is capitalized into prices. Credit for parents is the
exception: its effect on completed fertility is larger with a fixed stock
(7.8 against 6.3 percent), plausibly because the falling price makes ownership
cheaper, so access to ownership tied to a child is worth more. This reading is
inferred, not decomposed.

\paragraph{A more elastic supply.} With a supply elasticity of 1.5 from 2023 on,
a 10 percent shift in the supply curve lowers the price by 5.0 percent and
raises completed fertility by 2.6 percent; births per percentage point of
price decline are about the same as at $\eta=0.63$. Credit remains a matter of
timing at either elasticity.

\begin{table}[p]
\centering
\caption{Additional policies from 2023: effects relative to the no-policy path}
\label{tab:app-policy-more}
\small
\begin{tabularx}{\textwidth}{@{}Xrrrrr@{}}
\toprule
 & Price & Births & \multicolumn{2}{c}{Completed fertility} & Households \\
 \cmidrule(lr){4-5}
 & 2023 & 2023--26 & cohort 2023 & cohort 2043 & 2063 \\
\midrule
\multicolumn{6}{@{}l}{\emph{Estimated earnings risk}}\\
Property tax 1.06\% to 2\%, rebated & $-6.2$ & $-0.2$ & $+0.8$ & $+0.7$ & $+0.2$ \\
Earnings insurance & $+3.8$ & $+14.8$ & $+13.8$ & $+12.9$ & $+4.9$ \\
Flat transfer & $+2.0$ & $+8.0$ & $+7.8$ & $+7.0$ & $+2.7$ \\
Fixed stock (vs.\ elastic) & $-9.2$ & $-0.1$ & $+3.5$ & $+8.6$ & $+1.1$ \\
\quad Credit relaxation, LTV 100 & $+2.1$ & $+3.2$ & $+0.1$ & $-0.5$ & $+0.4$ \\
\quad Earnings insurance & $+12.8$ & $+13.0$ & $+10.1$ & $+7.8$ & $+3.9$ \\
\quad Flat transfer & $+7.0$ & $+7.1$ & $+5.5$ & $+3.9$ & $+2.1$ \\
\quad Credit for parents & $+3.5$ & $+7.9$ & $+7.8$ & $+6.3$ & $+2.7$ \\
Supply elasticity 1.5 (vs.\ 0.63) & $+0.6$ & $+1.0$ & $+0.4$ & $-1.4$ & $+0.1$ \\
\quad Housing supply $+10\%$ & $-5.0$ & $+1.5$ & $+2.6$ & $+2.4$ & $+0.8$ \\
\quad Credit relaxation, LTV 100$^{*}$ & $+0.3$ & $+2.1$ & $-0.7$ & $+0.5$ & $+0.2$ \\
\addlinespace
\multicolumn{6}{@{}l}{\emph{Lower earnings risk}}\\
Property tax 1.06\% to 2\%, rebated & $-6.4$ & $-0.7$ & $+0.5$ & $+0.7$ & $+0.1$ \\
Earnings insurance & $+2.6$ & $+7.4$ & $+6.4$ & $+6.3$ & $+2.6$ \\
Flat transfer & $+1.8$ & $+5.9$ & $+5.6$ & $+5.2$ & $+2.2$ \\
Fixed stock (vs.\ elastic) & $-2.7$ & $-0.9$ & $+0.5$ & $+3.0$ & $+0.2$ \\
\quad Credit relaxation, LTV 100 & $+2.0$ & $+3.1$ & $-0.9$ & $-0.3$ & $+0.2$ \\
\quad Earnings insurance & $+7.1$ & $+6.2$ & $+4.6$ & $+4.3$ & $+2.1$ \\
\quad Flat transfer & $+5.6$ & $+5.2$ & $+4.2$ & $+3.3$ & $+1.7$ \\
\quad Credit for parents$^{*}$ & $+4.9$ & $+11.8$ & $+10.5$ & $+8.5$ & $+4.2$ \\
Supply elasticity 1.5 (vs.\ 0.63) & $-0.3$ & $+0.6$ & $+0.5$ & $-0.3$ & $+0.2$ \\
\quad Housing supply $+10\%$ & $-5.1$ & $+1.3$ & $+2.1$ & $+1.9$ & $+0.7$ \\
\quad Credit relaxation, LTV 100 & $+0.7$ & $+3.7$ & $-0.2$ & $+0.1$ & $+0.3$ \\
\bottomrule
\end{tabularx}
\par\smallskip
\begin{minipage}{0.95\linewidth}
\footnotesize\textit{Notes:} Percent differences. Indented rows are compared
with the no-policy path under the same supply variant (fixed stock, or
elasticity 1.5); the unindented supply rows compare that no-policy path with
the baseline one.
\problem{Market clearing has not been established for the responses marked
$^{*}$.}
\end{minipage}
\end{table}

\subsection{The credit relaxation by cohort}
\label{app:policy-cohorts}

The credit relaxation lowers the completed fertility of the cohort aged 18 to
21 in 2023 by 1.1\%, raises it for the cohorts entering between the late 2030s
and the early 2060s, has almost no effect on the cohorts after them and raises
it again in the steady state (Figure~\ref{fig:app-cohort-2023}b, Baseline). To see
why, I solve both scenarios again with the paths of the house price (and with
it rents), the pension and the property-tax rebate held at their no-policy
values, so that only the borrowing limit differs. The effects on these paths
are close to the equilibrium ones (Table~\ref{tab:app-cohorts}), so they do
not come from the response of prices. On these paths the relaxation raises births in 2023
by 1.38\%, against 1.35\% in equilibrium. I then change one of these paths, the
same way in both scenarios.

\paragraph{The 2023 cohort.} The pension falls because the payroll tax is
fixed while smaller cohorts reach working age: retirees per worker rise from
0.35 in 2023 to 0.47 in 2071, when this cohort retires, and the pension falls
by a quarter (Figure~\ref{fig:app-cohort-2023}a). With the pension held at its
2023 level, the relaxation raises the cohort's completed fertility by 0.6\%.
Holding the rebate or the house price at its 2023 level instead leaves losses
of 0.5\% and 0.3\%. Households that buy earlier and with more debt under the
relaxation have to rebuild more wealth before a smaller pension: the fall in
the pension raises their net wealth at age 46 by 0.27, against 0.13 without
the relaxation, and they have about 2\% fewer children a year at ages 26 to 42.

\begin{figure}[H]
\centering
\includegraphics[width=0.95\linewidth]{\PRF/fig_appE_cohort_2023.pdf}
\caption{The credit relaxation by cohort, with one path held at its 2023 level}
\label{fig:app-cohort-2023}
\par\smallskip
\begin{minipage}{0.92\linewidth}
\footnotesize\textit{Notes:} Panel (a): house price, pension and retirees per
worker on the no-policy path, $2023=100$; the dotted line marks 2071, when the
cohort aged 18--21 in 2023 retires. Panel (b): \% difference in the completed
fertility of each cohort between the credit relaxation (loan-to-value 100) and
no policy. Baseline compares the two equilibrium paths. In the other lines both
scenarios face the no-policy paths of the house price, pension and rebate,
with the path named held at its 2023 value. Each point is one cohort.
\end{minipage}
\end{figure}

\paragraph{The later cohorts.} Cohorts entering after the early 2030s retire
once the pension has started to recover, and they gain from the relaxation
only while house prices keep falling. With the price held at its 2023 level,
the relaxation lowers the completed fertility of every cohort, by up to
1.4\%. With the price held at its 2063 level from 2063 on
(Figure~\ref{fig:app-cohort-dip}), the gain of the cohorts entering between
2039 and 2059 turns into a loss of 0.2--0.3\%, and the later cohorts are
within 0.1\% of zero. Without the stop, the cohorts entering between 2071 and
2095 live at prices between 0.50 and 0.60, where the steady-state effect of
the relaxation at a fixed price is also close to zero ($-0.11\%$ at 0.55,
$+0.05\%$ at 0.50, $+0.08\%$ at 0.61). For completed fertility, the second
fall in the effect is the end of the gain from falling prices, not a new loss.

\begin{figure}[H]
\centering
\includegraphics[width=0.95\linewidth]{\PRF/fig_appE_cohort_dip.pdf}
\caption{The credit relaxation by cohort, with the house price held from 2063}
\label{fig:app-cohort-dip}
\par\smallskip
\begin{minipage}{0.92\linewidth}
\footnotesize\textit{Notes:} Panel (a): house price on the no-policy path, and
the path held at its 2063 value from 2063 on. Panel (b): as in
Figure~\ref{fig:app-cohort-2023}b. In the second line both scenarios face the
no-policy pension and rebate and the held price path.
\end{minipage}
\end{figure}

Births in each period do not follow the same pattern, because a cohort's
completed fertility sums births over several periods. With the price held from
2063, the effect of the relaxation on births is $+0.8\%$ in 2047, $+0.1\%$ in
2055 and $-0.5\%$ in 2063, and between $0.1\%$ and $-0.2\%$ after that, so
the gain comes earlier and the fall around 2063 is deeper than in the
baseline. With the decline after 2055 at half its speed, the effect on births
stays between $+0.3\%$ and $+0.6\%$ from 2063 to 2103. The effect of the
relaxation on births follows how fast prices are expected to fall.

\begin{table}[H]
\centering
\caption{Credit relaxation: completed fertility by cohort}
\label{tab:app-cohorts}
\small
\input{\PRF/appE_cohort_table.tex}
\par\smallskip
\begin{minipage}{0.92\linewidth}
\footnotesize\textit{Notes:} \% difference in completed fertility between
the credit relaxation and no policy, by the year the cohort is aged 18--21.
Baseline: equilibrium paths. Other rows: both scenarios face the no-policy
paths of the house price, pension and rebate, with the changes named.
\end{minipage}
\end{table}

\subsection{How credit for parents works}
\label{app:policy-parents-mech}

Table~\ref{tab:app-parents} reports the numbers behind the description of
credit for parents in Section~\ref{sec:policy-parents}. Most young parents
already own without the policy, so ownership among households aged 22 to 29
with a child at home rises only from 75.6\% to 78.3\%. The policy works
through the decision to have a first child: childless renters aged 22 to 29
try for a first child 13 to 16\% more often in 2023. Once they have a child,
renters use the credit, buying with probability 0.63 instead of 0.57 at ages
22 to 25, and by 2031, 11\% of owners aged 22 to 45 with a child at home owe
more than 85\% of the value of their home. The credit relaxation for all
households, in the last column, instead raises the purchases of childless
renters, from 0.25 to 0.75 at ages 22 to 25, and in 2023 barely changes
their first-birth attempts.

\begin{table}[H]
\centering
\caption{Credit for parents and the credit relaxation, households in 2023 and 2031}
\label{tab:app-parents}
\small
\input{\PRF/appE_parents_table.tex}
\par\smallskip
\begin{minipage}{0.92\linewidth}
\footnotesize\textit{Notes:} Equilibrium paths from the same 2023 state.
Ownership is measured after the tenure choice of the period. Probabilities
are per four-year period: trying for a first child, and buying a home for a
household that starts the period as a renter. Debt relative to home value is
measured at the start of the period, among owners. Credit for parents raises
the loan-to-value limit to 1 for households with a child at home; the credit
relaxation raises it to 1 for all households.
\end{minipage}
\end{table}

\subsection{A descriptive check in the PSID}
\label{app:psid-small-home}

In the model, households that buy a small home before their first child are
less likely to have one. Table~\ref{tab:app-psid-small-home} compares women in
the PSID who bought their first home at ages 22 to 33. Among those who were
childless at purchase, 43\% of the buyers of a home of four rooms or fewer had
a first birth within six years, against 63\% of the buyers of larger homes;
with controls the gap is 14 percentage points (standard error 6.7). Among
buyers who already had one child, home size makes no difference to the second
birth. The comparison is descriptive: households that do not plan to have
children may also choose smaller homes.

\begin{table}[H]
\centering
\caption{First-time buyers in the PSID: births within six years by home size at purchase}
\label{tab:app-psid-small-home}
\small
\input{\PRF/appE_psid_table.tex}
\par\smallskip
\begin{minipage}{0.92\linewidth}
\footnotesize\textit{Notes:} PSID-SHELF, women who are household heads or
partners with a known birth history. First-time buyers own at an interview,
did not own at the previous interview (at most two years earlier) and never
owned before; they are aged 22 to 33 and bought in 1997--2001, 2003--2007 or
2009--2013, and are observed six years later. One observation per woman.
Rooms are counted at the first interview as an owner. The difference is from
a linear probability model with dummies for purchase period, three-year age
group, income decile, education (five groups) and living with a partner;
standard errors clustered by woman in parentheses. Unweighted.
\end{minipage}
\end{table}

\section{A CES Housing Composite}
\label{app:ces}

% Source identity retained: previous calibration loss 14.40, matched-state room target.
\problem{The CES comparison uses a different room-response target and
parameter vector, so it does not isolate the effect of the housing composite.}
The baseline combines a Cobb--Douglas composite with the space requirement
$h_P$. Here I use a CES composite in which children raise the rooms needed
for given housing services:
\begin{equation}
 u_t=\frac{1}{1-\sigma}\left(\frac{Q_t}{e(n_t)}\right)^{1-\sigma}+\psi_t n_t^{1-\gamma},\qquad
 Q_t=\Bigl[\alpha_0\,c_t^{\rho}+(1-\alpha_0)\Bigl(\frac{\chi^{\1\{\mathrm{own}\}}\,h_t}{k(n_t)}\Bigr)^{\rho}\Bigr]^{1/\rho},
\end{equation}
with $\rho=(\varepsilon-1)/\varepsilon$, $\varepsilon=0.67$, and
$k(n_t)=\exp\bigl[\eta_1\1\{n_t\ge1\}+\eta_2\max(n_t-1,0)\bigr]$. Housing and
consumption are complements. A child lowers the housing services each room
provides, and $\varepsilon$ governs how far consumption can make up for them.
Without a floor, space is not a fixed entry cost of a child.

I estimate $\eta_1$ and $\eta_2$, the rooms needed by the first and by each
later child at home, and $\alpha_0$, so twelve parameters match thirteen
moments. The new moments are the rent share of renters aged 30--55 with no
child at home (0.360, CE), the probability of moving from renting to owning
within four years at ages 22--41 (0.281, PSID), and the share of mothers with
three or more children (0.357, CPS). I re-measure the rooms added at the first
birth (1.025, PSID), target the room gap between larger and smaller families,
and use the ownership gap of recent parents as a validation moment.

% completed fertility at 46--49, CES at the calibrated point, fixed price and rebate (generated)
\def\cesLTVhundred{$+0.1\%$}
\def\cesLowerRisk{$+13.2\%$}
\def\cesCapTen{$+0.5\%$}

Table~\ref{tab:app-ces-fit} reports the fit. Against its matched
Cobb--Douglas comparison, the CES variant fits childlessness, family size,
rooms and the rooms at first birth slightly better, and timing, ownership,
the rent share, the family room gap, wealth and bequests worse. The responses
of completed fertility follow the pattern of the matched comparison:
\cesLTVhundred{} with 100\% financing, \cesLowerRisk{} with lower earnings
risk, and \cesCapTen{} with a ten-room rental cap.

\begin{table}[H]
\centering
\caption{CES variant: targets and fit}
\label{tab:app-ces-fit}
\small
\begin{minipage}[t]{0.66\linewidth}\vspace{0pt}
\begin{tabularx}{\textwidth}{@{}Xrrr@{}}
\toprule
Moment & Target & CES & Cobb--Douglas \\
\midrule
Childless, women 40--44 (\%) & 19.8 & 20.1 & 20.2 \\
One child among mothers, 40--44 (\%) & 21.4 & 21.5 & 21.7 \\
Mothers with 3+ children, 40--44 (\%) & 35.7 & 38.7 & 39.5 \\
Mean age at first birth (years) & 25.98 & 26.02 & 25.96 \\
Children ever born by age 25 & 0.810 & 0.542 & 0.553 \\
\addlinespace[3pt]
Mean occupied rooms & 5.73 & 5.80 & 5.82 \\
Ownership, heads 30--55 (\%) & 67.6 & 66.6 & 67.0 \\
Renter to owner within 4 years (\%) & 28.1 & 34.7 & 31.0 \\
Rent share, childless renters (\%) & 36.0 & 17.2 & 23.3 \\
First-birth room response & 1.025 & 0.801 & 1.264 \\
Room gap, 3+ vs 1--2 children at home & 0.385 & 0.538 & 0.294 \\
\addlinespace[3pt]
Wealth / annual gross earnings & 4.46 & 5.07 & 4.58 \\
Annual bequests / wealth (\%) & 0.73 & 0.86 & 0.71 \\
\bottomrule
\end{tabularx}

\end{minipage}\hfill
\begin{minipage}[t]{0.3\linewidth}\vspace{0pt}
\begin{tabular}{@{}lr@{}}
\toprule
Parameter & Value \\
\midrule
$\beta_{\rm annual}$ & 0.948 \\ % Annual discount factor
$\kappa_1$ & 0.100 \\ % First-birth taste scale
$\kappa_C$ & 0.246 \\ % Later-birth taste scale
$\chi$ & 1.061 \\ % Owner housing-service multiplier
$H_0$ & 7.941 \\ % Housing supply scale
$\theta_0$ & 0.421 \\ % Bequest-motive scale
$\xi$ & 0.304 \\ % First-birth utility cost
$\psi_0$ & 0.150 \\ % Initial preference for children
$\gamma$ & 0.156 \\ % Curvature of the benefit of children
$\kappa_T$ & 0.016 \\ % Tenure taste scale
$\alpha_0$ & 0.720 \\ % Consumption weight in the composite
$\eta_1$ & 0.630 \\ % Housing need of the first child
$\eta_2$ & 0.059 \\ % Housing need of each later child
\bottomrule
\end{tabular}

\end{minipage}
\end{table}

In Table~\ref{tab:app-ces-credit}, I vary only $\psi$ to match 19\%
childlessness under each specification, at fixed prices and with the
six-room rental cap. In all thirteen specifications that reach 19\%
childlessness, the credit relaxation moves completed fertility by between $-2\%$
and $+0.3\%$, and the negative sign comes mostly from the floor. Credit to
renters raises births on impact by 6\% to 9\% and completed fertility by
between $-0.6\%$ and $+2.4\%$ in the cases shown.

\begin{table}[H]
\centering
\caption{CES variant: credit across specifications}
\label{tab:app-ces-credit}
\small\setlength{\tabcolsep}{4pt}
\begin{tabularx}{\textwidth}{@{}Xccc rr rr@{}}
\toprule
 & & & & \multicolumn{2}{c}{Financed share $80\%\to95\%$} & \multicolumn{2}{c}{Renter credit, 2 years} \\
\cmidrule(lr){5-6}\cmidrule(l){7-8}
 & $\varepsilon$ & Floor & Need $k$ & Impact & Completed & Impact & Completed \\
\midrule
Cobb--Douglas (benchmark) & 1 & yes & 1 & $+0.42\%$ & $-1.35\%$ & $+8.5\%$ & $+0.62\%$ \\
CES & 0.5 & yes & 1 & $-0.16\%$ & $-1.40\%$ & $+8.0\%$ & $+0.22\%$ \\
CES & 0.25 & yes & 1 & $-0.06\%$ & $-0.87\%$ & $+7.3\%$ & $-0.25\%$ \\
CES & 0.5 & yes & 1.5 & $+0.05\%$ & $-1.10\%$ & $+7.0\%$ & $-0.63\%$ \\
CES, no floor & 0.5 & no & 1 & $+0.44\%$ & $-0.03\%$ & $+6.9\%$ & $+2.25\%$ \\
CES, no floor & 0.5 & no & 1.5 & $+0.51\%$ & $-0.18\%$ & $+8.3\%$ & $+2.39\%$ \\
Cobb--Douglas, no floor & 1 & no & 1.5 & $+1.03\%$ & $+0.23\%$ & $+8.3\%$ & $+2.42\%$ \\
\bottomrule
\end{tabularx}
\par\smallskip
\begin{minipage}{0.95\linewidth}
\footnotesize\textit{Notes:} Each row varies only $\psi$ to match 19\%
childlessness. House prices are fixed and the rental cap is six rooms.
\end{minipage}
\end{table}

Completed fertility moves little with the general credit relaxation in every
specification, so the timing response does not hinge on the Cobb--Douglas
composite. These comparisons hold prices fixed and use other parameters and
targets, so they do not measure the contribution of the floor to the results
of Section~\ref{sec:policy}.

\section{Solution Algorithms}
\label{app:algorithms}

\paragraph{Stationary equilibrium.} I solve for the house price $P$ and the
rebate $T$, given the parameters, the preference for children and the payroll
tax. In a stationary population the age distribution follows from survival
and entry, so the pension that balances the payroll tax is known before the
search. I check the pension budget at the end.
\begin{enumerate}
\item Guess $P$ and compute the rent $r$ from the user cost.
\item Solve the household problem backward over age and compute the choice
probabilities.
\item Compute the stationary distribution of households, iterating on the
rebate until $T$ times the number of households equals property-tax revenue.
In the calibration, the number of households is one and $H_0$ is set so that
supply equals demand at $P$. In the counterfactuals, $H_0$ is fixed and the
number of households equals housing supply divided by mean housing per
household.
\item Compute the renewal residual: adjusted births divided by 2.1, relative
to the number of entrants, minus one. It is zero when the population
reproduces itself.
\item Update $P$ by bracketing the root of the renewal residual, and stop when
the renewal and rebate residuals are below $10^{-6}$. Re-solve at the
selected price, require the same residuals, and check that the pension budget
balances to the same tolerance.
\end{enumerate}
% [CODE] native_phase_b.py 25-26 (RENEWAL_TOL, PAYGO_TOL 1e-6), 97-99
% (supply rule), 126-127 (population = supply/demand; renewal = births/(2.1
% entry) - 1), 440-454 (bracket, selected root, exact repeat);
% equilibrium.py (model) 84-96 (implied H0 under population_one, renewal
% root); fiscal_closure.py (T N = revenue - grants, FISCAL_TOL 1e-6);
% e5f_stationary_paygo.py 47-62 (pension from the payroll tax).

\paragraph{Transition.} I solve for the paths of the price, the pension and
the rebate, $(P_t,m_t,T_t)_{t\geq t_0}$, given the inherited households and
the path of $\psi_t$, which households foresee.
\begin{enumerate}
\item Solve the terminal stationary equilibrium. It gives the terminal value
function.
\item Guess the three paths. Compute rents from \eqref{eq:m-rent} and, with a
slowly adjusting stock, the stock from \eqref{eq:app-slow-stock}.
\item Solve the household problem backward in calendar time from the terminal
value.
\item Move households and the queue of entrants forward from the inherited
state, and add up housing demand and taxes at each date.
\item Compute the residuals of the housing market, the pension budget and the
rebate at every date. The loss is the largest residual across dates.
\item Update the three paths jointly and repeat until the loss is below
tolerance. Extend the horizon to check that early responses are stable.
\end{enumerate}
With free entry and $\eta=0.63$, the transitions use 100 four-year periods,
from 2007 (2023 for the policy experiments) to 2403. With $\eta=1.5$ and with
a slowly adjusting stock, they use 200 periods, to 2803. The tolerances are
$2\times10^{-4}$ on the housing market, $2\times10^{-5}$ on the pension and
the rebate, and $10^{-6}$ on the renewal condition of the terminal stationary
equilibrium.
% [MODEL, Oct 24] Horizons and residuals read from the result files: free entry eta 0.63
% 20261022_fi_all_arms/history/FIN063fi_fw.json (100 periods) and data_fi_torch/P063fi_*.json
% (96 periods from 2023); eta 1.5 FIN15fi200_w.json and P15fi_*.json (200/196); slow stock
% data_k200_torch/FIN063k200_f_history.json and P063k_*.json (200/196). Converged arms: max
% market residual <= 2.0e-4, fiscal and rebate <= 1.9e-5.
